% !TeX program = xelatex
% 东三省/辽宁省数学建模竞赛模板 — 基于 nemcmthesis 文档类
\documentclass{nemcmthesis}

% === 竞赛信息 ===
\ttle{[论文题目]}
\category{[研究生/本科生/专科生]}
\school{[学校名称]}
\tihao{[A/B/C/D]}
\membera{[队员1]}
\memberb{[队员2]}
\memberc{[队员3]}
\phonea{[电话1]}
\phoneb{[电话2]}
\phonec{[电话3]}
\supervisor{[指导教师]}

% === 额外宏包 ===
\usepackage{tikz}
\usetikzlibrary{arrows.meta,positioning,shapes.geometric,fit,calc}
% ⛔ algorithm2e 已由 nemcmthesis.cls 加载，不要重复加载
\usepackage[section]{placeins}
\usepackage{needspace}
\usepackage{longtable}

% === 引用上标格式（nemcmthesis 用 natbib，加 super 选项）===
\setcitestyle{super}

\begin{document}

\makecoverpage

% === 摘要 ===
\input{sections/abstract}

% === 正文 ===
\input{sections/1_restatement}
\input{sections/2_analysis}
\input{sections/3_assumptions}
\input{sections/4_symbols}
\input{sections/5_problem1}
\input{sections/6_problem2}
\input{sections/7_problem3}
\input{sections/8_evaluation}

% === 参考文献 ===
\begin{thebibliography}{99}
% \bibitem{ref1} 作者. 标题[J]. 期刊, 年份.
\end{thebibliography}

% === 附录 ===
\input{sections/A_code}

\end{document}
